% s-040-p040-m003; original PDF 40, 41
\Aphorism{s-040-p040-m003}{26}
Auch der Philosoph nach Nietzsches Traum, der cäsarische Züchter und Gewaltherr mit ökumenischen Zielen, ist im Grunde ein recht unbedenklicher Altruist. Ob man die „Andern“ erfreuen, bessern, drücken, lenken, züchten will: das Gemeinsame aller dieser Versuche ist, dass sie ein Problem von der irrationalen Seite anfassen, da wo entweder der Erfolg eines gegebenen Mittels oder das Mittel zu gegebenem Zweck unbekannt ist oder Beides. Der Egoist greift hingegen da zu, wo eine Handhabe ist — der Egoist im weitesten Sinne, also nicht nur Der, der auf seinen Vortheil, sondern ebenso der auf seine Erziehung, seine Erkenntniss, unter Umständen seine Selbstpreisgebung aus ist; ich sehe zwischen „Selbstsucht“ und Selbstlosigkeit keinen Unterschied des Zieles, sondern nur der Methode, und finde alle Vernunft und Wissenschaft auf Seite des Egoisten. Wie man sich selbst ein Stück in gegebener Richtung vorbringt, weiss man zur Noth; beim Andern läuft man stets Gefahr, Nichts oder das Gegentheil des Gewollten zu erreichen. Mit der „Erlösung“ des Egoismus zur allgemeinen Menschenliebe vergrössert man nicht die Moralität, sondern das Gebiet, in dem man Dummheiten anrichten kann
\par

